% Generated draft: requires mathematical and source review.
\documentclass[11pt]{article}
\usepackage[margin=25mm]{geometry}
\usepackage{fontspec}
\setmainfont{Latin Modern Roman}
\usepackage{amsmath,amssymb,mathtools,unicode-math}
\setmathfont{Latin Modern Math}
\usepackage{enumitem,needspace,xurl,hyperref}
\setlength{\parindent}{0pt}
\setlength{\parskip}{5pt}
\setlength{\emergencystretch}{3em}
\newcommand{\N}{\mathbb{N}}
\newcommand{\Z}{\mathbb{Z}}
\newcommand{\Q}{\mathbb{Q}}
\newcommand{\R}{\mathbb{R}}
\newcommand{\C}{\mathbb{C}}
\newcommand{\F}{\mathbb{F}}
\newcommand{\A}{\mathbb{A}}
\newcommand{\PP}{\mathbb{P}}
\newcommand{\E}{\mathbb{E}}
\newcommand{\eps}{\varepsilon}
\newcommand{\dd}{\,\mathrm{d}}
\begin{document}
\section*{UnsolvedMath ID: 30002099}
\textbf{Source problem code:} catalogue entry
\subsection{Definitions and mathematical statement}
Let \(k\) be a number field and let \(X/k\) be a smooth projective surface. Under the usual interpretation, a regular polarized self-map is an everywhere-defined morphism \(f\colon X\to X\) for which there exist an ample line bundle \(L\) on \(X\) and an integer \(q>1\) such that \(f^*L\simeq L^{\otimes q}\). The supplied record does not explicitly state this polarization convention.

For a geometric periodic point \(p\in X(\overline{k})\), let \(n_p\geq 1\) be its least period, so that \(f^{n_p}(p)=p\). The derivative of its return map is
\[
\mathrm{d}(f^{n_p})_p\colon T_pX_{\overline{k}}\longrightarrow T_pX_{\overline{k}},
\]
with two eigenvalues \(\lambda_1(p),\lambda_2(p)\in\overline{k}\), counted with multiplicity. Two nonzero eigenvalues are multiplicatively independent if
\[
\forall(a,b)\in\Z^2,\qquad
\lambda_1(p)^a\lambda_2(p)^b=1
\ \Longrightarrow\ (a,b)=(0,0).
\]
Under the geometric-point and return-map interpretation, the exceptional condition in the question is
\[
\nexists p\in X(\overline{k})\quad
\left[
 p\text{ is periodic},\quad
 \lambda_1(p)\lambda_2(p)\ne 0,\quad
 \lambda_1(p),\lambda_2(p)\text{ are multiplicatively independent}
\right].
\]
This is a conditional normalization: the source does not specify the field of definition of the periodic points or explicitly identify the derivative as that of the return map. Multiplicative independence in the displayed sense applies only to nonzero eigenvalues.

\textbf{Source question, with unresolved qualifications.} Classify the regular polarized self-maps of smooth projective surfaces over number fields satisfying the exceptional condition, and determine whether the invariant algebraic curves intended by the source admit a natural parametrization by elements of \(\overline{\Q}\) rather than \(\Q\).

The record does not specify the equivalence relation for the classification, which curves are meant by the resulting invariant curves, or the mathematical meaning of the proposed parametrization. In particular, it does not determine whether a curve \(C\subset X_{\overline{k}}\) must satisfy \(f(C)=C\), whether \(f^r(C)=C\) for some \(r\geq1\) suffices, or whether irreducibility is required. Nor does it distinguish an indexing of curves, a parameter space for a family, or a parametrization of each individual curve. Accordingly, these definitions isolate an interpretable dynamical condition but do not yield a uniquely specified classification or parametrization problem. The supplied literature assessment also refers to an unidentified conjecture of Zhang; its statement and role cannot be recovered from this record.
\subsection{Short English statement}
The question asks for a classification of polarized surface self-maps having no periodic point whose return-map derivative has two multiplicatively independent nonzero eigenvalues, and for a natural description of certain associated invariant curves using \(\overline{\Q}\) rather than \(\Q\). The periodic-point domain, the intended curves, and the meaning of parametrization remain unspecified, so this is a qualified statement-only draft.
\subsection{Sources}
\begin{itemize}
\item Supplied local UnsolvedMath registry record, permanent ID \(30002099\), title: \emph{Exceptional Dynamics of Polarized Surface Endomorphisms}; collection: \emph{Oberwolfach Reports Open Problems}; exact source problem code: \texttt{catalogue entry}. No individual proposer is named; the supplied proposed year is 2012.
\item Supplied report citation: \emph{Classical Algebraic Geometry (2013)}, Oberwolfach Reports, pp.\ 1845--1893, DOI \texttt{10.4171/owr/2012/30}. Supplied source URL and literature-source link, labeled ``original OWR report'': \url{https://doi.org/10.4171/owr/2012/30}. The record describes these reports as documentation of mathematical workshops and their open-problem discussions.
\item Supplied literature assessment, dated exactly \texttt{2026-08-21}: status labeled open; the extracted question does not identify Zhang's conjecture or the intended class of invariant curves sufficiently to test a theorem, and therefore remains unresolved under the supplied normalization. This is supplied dated curation derived from the OWR record, not an independently verified literature assessment. Its instruction to verify time-sensitive status and cited claims independently remains applicable.
\item The registry labels the statement \texttt{corrected\_verified} and reports that its mathematical content was checked against the official report, with notation, context, and question form normalized without changing the source-level claim. These are supplied verification claims; no source visit or independent mathematical verification is asserted here.
\end{itemize}
\end{document}
